\chapter{Introduction}

The field of bio-inspired robotics \cite{delcomyn2004insect} has seen rapid advancements in recent years,
\cite{chai2022survey,fan2024review} driven by the need for autonomous systems capable of navigating
unstructured and complex environments. Among these, hexapod robots offer a unique combination of stability,
payload capacity, and terrain adaptability that wheeled platforms cannot match. However, the effective deployment
of such legged systems is fundamentally dictated by the quality of their actuation.


\subsubsection{Motivation and Problem Statement}

Precise legged locomotion requires accurate control of joint angles and torque. While industrial-grade actuators
(e.g. Dynamixel-X series \cite{robotis_dynamixel_x}) offer the necessary fidelity and communication interfaces,
their prohibitive cost makes them inaccessible for many research and prototyping applications. Conversely,
low-cost "hobby" servo-motors, while affordable, present critical limitations. They typically operate as
"black boxes" with inaccessible internal control loops, rely on wear-prone resistive potentiometers for feedback,
and suffer from significant deadbands and non-linear friction effects.


\subsubsection{Thesis Objectives}

This thesis addresses these challenges by proposing a comprehensive hardware-software design. The overarching goal
is to bridge the gap between low-cost hardware and high-performance control. The specific objectives are:

\begin{enumerate}
    \item \textbf{Kinematic Framework:} To derive and validate the mathematical model of a 3-DOF hexapod limb,
        ensuring reliable workspace analysis and trajectory planning.
    \item \textbf{Actuator Redesign:} To engineer a custom "Smart Actuator" module that retains the standard form factor
        of commercial servos but upgrades performances through internal electronics replacements.
    \item \textbf{Sim-to-Real Validation:} To establish a Hardware-in-the-Loop (HIL) architecture that allows for direct
        performance comparison between a theoretical numerical model and the physical prototype.
    \item \textbf{Modular Scalability:} To adopt a strictly modular design philosophy for both hardware and software
        architectures, ensuring that the single-leg system can be seamlessly replicated and integrated to realize the
        complete hexapod robot in future developments.
\end{enumerate}


\subsubsection{Thesis Structure}

The remainder of this thesis is organized as follows:

\begin{itemize}
    \item Chapter \ref{ch:model} establishes the theoretical foundation. It details the mechanical design of the limb,
        the Denavit--Hartenberg parameterization, and the Inverse Kinematics derivation. It also presents the
        numerical model developed in Simscape Multibody, which serves as reference for validation.
    \item Chapter \ref{ch:servomotor} focuses on the component-level engineering. It analyzes the limitations of
        commercial servos and details the hardware and software design of the custom actuator module, including the
        Switching PID control algorithm, ultimately evaluating the achieved improvements.
    \item Chapter \ref{ch:prototype} describes the system integration and experimental validation. It outlines the
        Hardware-in-the-Loop architecture and presents a quantitative comparison of the trajectory tracking performance
        between the simulation and the physical prototype.
    \item Chapter \ref{ch:conclusions} concludes the thesis by summarizing the achieved results, and outlining future
        research directions.
\end{itemize}

